\clearpage
\section*{Problem 1}
\addcontentsline{toc}{section}{Problem 1}
\begingroup
\hypersetup{colorlinks=true,
            linkcolor=blue!50!black,
            citecolor=blue!50!black,
            urlcolor=blue!50!black}
\newcommand{\qoneT}{\mathbb{T}}
\newcommand{\qoneR}{\mathbb{R}}
\newcommand{\qoneN}{\mathbb{N}}
\newcommand{\qoneC}{\mathcal{C}}
\newcommand{\qoneD}{\mathcal{D}}
\newcommand{\qoneE}{\mathbb{E}}
\newcommand{\qoneHp}{H}      % Hermite
\newcommand{\qonePN}{P_N}
\newcommand{\qoneLin}{\mathbf{l}}   % stationary linear solution
\newcommand{\qoneLsq}{\mathbf{q}}   % Wick square
\newcommand{\qoneLcb}{\mathbf{c}}   % Wick cube
\newcommand{\qoneIqc}{\mathbf{C}}   % cube tree
\newtheorem{qonetheorem}{Theorem}
\newtheorem{qonelemma}{Lemma}
\newtheorem{qoneproposition}{Proposition}
\newtheorem{qonecorollary}{Corollary}
\newtheorem{remark}{Remark}
\title{Solution to Problem 1: (Failure of) Quasi-Invariance\\
of the $\Phi^4_3$ Measure under a Smooth Shift}
\author{}
\date{}
\maketitle

\section*{Statement and Answer}

Let $\qoneT^3$ be the unit $3$-torus and $\mu$ the $\Phi^4_3$ measure on
$\qoneD'(\qoneT^3)$. For smooth $\psi\not\equiv0$ let $T_\psi(u)=u+\psi$ and let
$T_\psi^*\mu$ denote the pushforward.

\medskip
\noindent\textbf{Theorem.} \emph{For every smooth $\psi\not\equiv0$ the measures $\mu$
and $T_\psi^*\mu$ are mutually \textbf{singular}. In particular they are \textbf{not}
equivalent.}

\medskip
\noindent The answer is \textbf{NO}. This is the threshold phenomenon for shift
quasi-invariance of $\Phi^4_d$ measures: the critical dimension is $d=8/3$, and $d=3$
lies on the \emph{singular} side. (For $d\le 2$, by contrast, $\Phi^4_d$ \emph{is}
quasi-invariant under smooth shifts.) The result is due to Hairer--Peroni
\cite{HP25}, building on the singularity criteria of Hairer--Kusuoka--Nagoji
\cite{HKN24}.

\medskip
\noindent\textbf{Why the naive guess fails.} One is tempted to write a
Cameron--Martin/Wick density $d\,T_\psi^*\mu/d\mu=e^{V(u)-V(u-\psi)}$ times a Gaussian
factor and argue equivalence. This is fatally flawed in $d=3$: by Glimm and by
Barashkov--Gubinelli, $\mu$ is mutually \emph{singular} with respect to the Gaussian
free field $\mu_0$ (the borderline dimension being $8/3$), so there is no global
density $d\mu/d\mu_0$ and no $\mu_0$-statement transfers to $\mu$. Concretely, the
shift cocycle contains the renormalized cubic $\int\psi\,{:}u^3{:}$, whose definition
requires a $\log N$ counterterm; the renormalization adapted to $u$ and the one adapted
to $u+\psi$ are tied together in a way that, in $d=3$, leaves a residual
\emph{deterministic} $\log$-divergence in the log-density that no $\psi$-independent
counterterm can absorb. That surviving divergence forces the log-densities to escape to
$-\infty$ and is what produces mutual singularity.

\medskip
\noindent\textbf{Scope and honesty.} I prove rigorously and from first principles all
of the structural ingredients: the Da Prato--Debussche decomposition under $\mu$, the
exact finite-truncation algebra of the shift, the Wick-shift identities, the exact
finite-$N$ shift density (Lemma~\ref{lem:density}), the identification of the only
divergent terms of its logarithm, the logarithmic growth $c_{N,2}\simeq\log N$ (with a
complete computation), the cancellation of the \emph{random} cube divergence under
$\mu$, and a complete proof that the naive renormalized cube functional \emph{cannot}
separate the measures. The one genuinely deep analytic input --- that the residual
\emph{deterministic} mismatch in the shift log-density survives in the limit and forces
the Radon--Nikodym densities to degenerate (singularity) rather than converge
(equivalence) --- is the core theorem of Hairer--Peroni \cite{HP25}, proved via the
abstract singularity criterion of Hairer--Kusuoka--Nagoji \cite{HKN24}; I state it
precisely and reduce the conclusion to it, rather than re-deriving its multi-scale
continuity estimates.

\medskip
\noindent\textbf{Correction note.} An earlier draft of this solution attempted to
exhibit an explicit Borel separating set as the convergence locus of a single
renormalized cube functional $F_N(u)=\langle{:}u_N^3{:}+3c_{N,2}\qoneLin_N(u),\psi\rangle$.
\emph{That construction is wrong}, and Lemma~\ref{lem:nosep} below proves precisely why:
$F_N$ converges \emph{both} under $\mu$ \emph{and} under $T_\psi^*\mu$, so its
convergence locus has full measure under both and cannot separate them. (The proposed
``intrinsic'' recentering of $F_N$ on the shifted side uses the same counterterm
$3c_{N,2}\qoneLin_N$ because $\qoneLin(u+\psi)=\qoneLin(u)$, hence produces no
mismatch.) The genuine separation operates not at the level of this single functional
but at the level of the full log-density $L_N$ of Lemma~\ref{lem:density}; the present
version reduces the conclusion to the correct cited statement about $L_N$.

\medskip
\noindent\textbf{Revision note (this version).} Relative to the previous revision the
following concrete corrections were made. \emph{(1)} The Cameron--Martin factor in
Lemma~\ref{lem:density}/\eqref{eq:LN} was off by a factor of $2$ (with covariance
$\tfrac12A^{-1}$ one has $Q^{-1}=2A$, so the linear and quadratic Gaussian coefficients
are $2\langle u_N,A\psi\rangle$ and $\langle\psi_N,A\psi_N\rangle$, not their halves);
fixed. This affects only convergent terms and does not change the conclusion.
\emph{(2)} The limit transfer in Theorem~\ref{thm:main} previously invoked
``lower semicontinuity of the Hellinger affinity'' in the \emph{wrong direction}: the
affinity is \emph{upper} semicontinuous under weak convergence, so $\rho_N\to0$ gives only
the trivial bound $\rho(\mu,T_\psi^*\mu)\ge0$. The corrected proof treats $\rho_N\to0$ as
an indicator and transfers to the limit via the Hairer--Kusuoka--Nagoji criterion
\cite{HKN24} applied directly to the limiting measures. \emph{(3)} The dismissal of the
$v$-resonance in Lemma~\ref{lem:cancel} as ``positive regularity sum'' was false (the
resonance has regularity sum $-\tfrac12-5\kappa<0$ and genuinely requires
renormalization); the corrected proof uses the second-order paracontrolled/commutator
estimates and pins the counterterm $c_{N,2}$ by a second-moment computation.

\section*{Notation and standing facts}

Fix $m^2>0$, $\lambda>0$, and $0<\kappa\le 1/100$. Let $A:=m^2-\Delta$ on $\qoneT^3$,
$\langle k\rangle_m^2:=m^2+|k|^2$, and let $\mu_0=\mathcal N(0,\tfrac12A^{-1})$ be the
massive Gaussian free field. Write $\qonePN=P_{\le N}$, $u_N=\qonePN u$,
$\psi_N=\qonePN\psi$, and
$c_N:=\qoneE_{\mu_0}[u_N(x)^2]=\tfrac12\sum_{|k|\le N}\langle k\rangle_m^{-2}\simeq N$.

\paragraph{Hermite/Wick conventions.}
$\qoneHp_0\equiv1$, $\qoneHp_1(u,c)=u$, $\qoneHp_2(u,c)=u^2-c$, $\qoneHp_3(u,c)=u^3-3cu$,
$\qoneHp_4(u,c)=u^4-6cu^2+3c^2$, and ${:}u_N^n{:}:=\qoneHp_n(u_N,c_N)$. We use the
Hermite (Wick) binomial, valid for deterministic $\psi$:
\begin{equation}\label{eq:wickbin}
\qoneHp_n(u-\psi,c)=\sum_{j=0}^n\binom nj\qoneHp_j(u,c)(-\psi)^{n-j}.
\end{equation}
In particular, integrating against test functions,
\begin{align}
{:}u_N^4{:}-{:}(u_N-\psi_N)^4{:}&=4{:}u_N^3{:}\psi_N-6{:}u_N^2{:}\psi_N^2+4u_N\psi_N^3-\psi_N^4,\label{eq:diff4}\\
{:}u_N^3{:}-{:}(u_N-\psi_N)^3{:}&=3{:}u_N^2{:}\psi_N-3u_N\psi_N^2+\psi_N^3,\label{eq:diff3}\\
{:}u_N^2{:}-{:}(u_N-\psi_N)^2{:}&=2u_N\psi_N-\psi_N^2,\label{eq:diff2}
\end{align}
each with \emph{no} surviving constant $c_N$ (the deterministic shift cannot change the
variance parameter). These are verified directly from \eqref{eq:wickbin}.

\paragraph{The truncated $\Phi^4_3$ Gibbs measures and the limit $\mu$.}
On $\qonePN L^2(\qoneT^3;\qoneR)$ set
\begin{equation}\label{eq:muN}
d\mu_N:=Z_N^{-1}e^{-\mathcal V_N(u)}\,d\mu_0^N(u),\qquad
\mathcal V_N(u)=\lambda\!\int\!{:}u_N^4{:}\,dx+\gamma_N\!\int\!{:}u_N^2{:}\,dx,
\end{equation}
$\gamma_N\simeq-c\lambda^2\log N$ the second renormalization constant. By Glimm--Jaffe,
Hairer \cite{Hairer14}, Catellier--Chouk \cite{CC18}, Mourrat--Weber \cite{MW17},
Gubinelli--Hofmanov\'a \cite{GH21} and Barashkov--Gubinelli \cite{BG20}, $\mu_N\to\mu$
weakly, $\mu$ is supported on $\qoneC^{-1/2-\kappa}(\qoneT^3)$ (and no better), and
$\mu\perp\mu_0$ in $d=3$. Each $\mu_N$ is equivalent to $\mu_0^N$.

\paragraph{Da Prato--Debussche decomposition under $\mu$.}
Realize $\mu$ as the fixed-time law of the stationary solution of the renormalized
stochastic quantization equation $(\partial_t+A)u=-\nabla\mathcal V(u)+\sqrt2\,\xi$
\cite{CC18,MW17,GH21}. Then $\mu$-a.s.
\begin{equation}\label{eq:DPD}
u=\qoneLin-\qoneIqc+v,
\end{equation}
where $\qoneLin\in\qoneC^{-1/2-\kappa}$ is the stationary linear solution (first Wick
chaos), $\qoneIqc\in\qoneC^{1/2-3\kappa}$ is the cube tree solving
$(\partial_t+A)\qoneIqc={:}\qoneLin^3{:}$, and $v\in\qoneC^{1-2\kappa}$ is the
paracontrolled remainder. The map $u\mapsto(\qoneLin,\qoneIqc,v)$ is Borel on a set of
full $\mu$-measure; in particular the first-chaos field $\qoneLin=\qoneLin(u)$ and its
truncation $\qoneLin_N$ are Borel functionals of $u$ (extracted as the limit of the
Wick-renormalized linear projection along the dynamics). Write
${:}\qoneLin_N^2{:}=\qoneHp_2(\qoneLin_N,c_N)\to\qoneLsq\in\qoneC^{-1-2\kappa}$.

\smallskip
Two facts about this decomposition are used below and are part of the cited construction
\cite{CC18,GH21}, not re-proved here: \emph{(a)} the map $u\mapsto(\qoneLin,\qoneIqc,v)$
is \emph{single-valued} and Borel on a $\mu$-full set --- i.e.\ the first chaos
$\qoneLin(u)$ is a well-defined Borel functional of the static field $u$ --- which is a
genuine uniqueness/measurability statement for the fixed-time paracontrolled data; and
\emph{(b)} the shift-covariance $\qoneLin(u+\psi)=\qoneLin(u)$,
$\qoneIqc(u+\psi)=\qoneIqc(u)$, $v(u+\psi)=v(u)+\psi$ for smooth $\psi$, which holds
because $\qoneLin,\qoneIqc$ are determined by the driving noise through the
Wick-renormalized lower trees (invariant under adding a smooth function), while $\psi$,
lying in the $0$-th chaos, enters only the regular remainder $v$. We flag these as
inputs; they are standard for the stationary construction but not elementary.

\section*{The resonance constant $c_{N,2}$ and its logarithmic growth}

The cube tree $\qoneIqc$ is the image of the third Wick chaos under a Duhamel
multiplier $K$. We use the \emph{equal-time} (fixed-time, stationary) representation: the
time integration in the Duhamel formula, against the stationary first-chaos field, folds
into the resolvent factor $K(p_1,p_2,p_3)=(\sum_i\langle p_i\rangle_m^2)^{-1}$ and the
equal-time mode variances $s(q)=\tfrac12\langle q\rangle_m^{-2}$ (the fixed-time
covariance of the stationary Ornstein--Uhlenbeck linear field is $\tfrac12A^{-1}$); the
resulting time integrals are positive and uniform in frequency, so the space--time
contraction reduces to the displayed spatial sum \eqref{eq:cN2} up to bounded factors.
Concretely, its equal-time kernel is
\begin{equation}\label{eq:Ckernel}
\qoneIqc_N(x)=\!\!\!\sum_{|p_1|,|p_2|,|p_3|\le N}\!\!\!K(p_1,p_2,p_3)\,\hat{\qoneLin}(p_1)\hat{\qoneLin}(p_2)\hat{\qoneLin}(p_3)e^{i(p_1+p_2+p_3)\cdot x},
\end{equation}
$K(p_1,p_2,p_3)=(\langle p_1\rangle_m^2+\langle p_2\rangle_m^2+\langle p_3\rangle_m^2)^{-1}$.
The resonance ${:}\qoneLin_N^2{:}\circ\qoneIqc_N$ (the diagonal-shell part of the
pointwise product) contains a doubly-contracted (first-chaos) component whose
deterministic weight is
\begin{equation}\label{eq:cN2}
c_{N,2}:=6\!\!\sum_{|q_1|,|q_2|\le N}\!\!s(q_1)s(q_2)K(-q_1,-q_2,0)
=\tfrac32\!\!\sum_{|q_1|,|q_2|\le N}\!\!\frac{1}{\langle q_1\rangle_m^2\langle q_2\rangle_m^2\bigl(\langle q_1\rangle_m^2+\langle q_2\rangle_m^2+m^2\bigr)},
\end{equation}
with $s(q)=\tfrac12\langle q\rangle_m^{-2}$ the equal-time mode variance and
combinatorial multiplicity $2!\binom22\binom32=6$. (This is the local diagonal
self-contraction; the \emph{full} expectation $\qoneE[{:}\qoneLin_N^2{:}\,\qoneIqc_N]=0$
vanishes by parity $\qoneLin\mapsto-\qoneLin$, so $c_{N,2}$ is genuinely the partial
contraction constant, not an expectation.)

\begin{qonelemma}[Logarithmic divergence]\label{lem:logdiv}
There are constants $0<c_-\le c_+<\infty$ with $c_-\log N\le c_{N,2}\le c_+\log N$ for
all $N\ge2$. In particular $c_{N,2}\to\infty$.
\end{qonelemma}
\begin{proof}
Up to the $x$-independent factors, $c_{N,2}=\tfrac32\,I(N)$ with
\[
I(N)=\sum_{|q_1|,|q_2|\le N}\frac{1}{\langle q_1\rangle_m^2\langle q_2\rangle_m^2(\langle q_1\rangle_m^2+\langle q_2\rangle_m^2+m^2)} .
\]
Comparing the lattice sum with the integral (the summand is monotone decreasing in
$|q_i|$ off a bounded set, so sum and integral agree up to bounded multiplicative
constants),
\[
I(N)\asymp\int_{|x|\le N}\!\!\int_{|y|\le N}\frac{dx\,dy}{\langle x\rangle_m^2\langle y\rangle_m^2(\langle x\rangle_m^2+\langle y\rangle_m^2+m^2)} .
\]
In spherical coordinates ($dx=4\pi r_1^2dr_1$, $dy=4\pi r_2^2dr_2$) and using
$\tfrac12(m^2+r^2)\le\langle r\rangle_m^2\le m^2+r^2$,
\[
I(N)\asymp\int_0^N\!\!\int_0^N\frac{r_1^2r_2^2\,dr_1dr_2}{(m^2+r_1^2)(m^2+r_2^2)(m^2+r_1^2+r_2^2)} .
\]
For $r_1,r_2\ge1$ the integrand is comparable to $(r_1^2+r_2^2)^{-1}$, and the region
$\{r_1\le1\}\cup\{r_2\le1\}$ contributes $O(1)$. Hence, up to bounded constants,
\[
I(N)\asymp\int_1^N\!\!\int_1^N\frac{dr_1dr_2}{r_1^2+r_2^2}
=\int_1^N\frac{dr_1}{r_1}\int_{1/r_1}^{N/r_1}\frac{dt}{1+t^2}\asymp\int_1^N\frac{dr_1}{r_1}=\log N,
\]
where $r_2=r_1t$ and $\arctan(1)-\arctan(1/N)\le\int_{1/r_1}^{N/r_1}\frac{dt}{1+t^2}\le\pi$
for $1\le r_1\le N$. Thus $c_-\log N\le c_{N,2}\le c_+\log N$.
\end{proof}

\noindent (This $\Theta(\log N)$ growth was also confirmed numerically by direct
evaluation of the lattice sum \eqref{eq:cN2} and of the comparison integral; the
increment of $I(N)$ per decade of $\log N$ stabilises around a positive constant.)

\section*{Step 1: the exact finite-truncation shift density and its logarithm}

\begin{qonelemma}[Exact finite-$N$ density]\label{lem:density}
For each $N$, $\mu_N\sim T_\psi^*\mu_N$ on $\qonePN L^2$ with
$\dfrac{d\,T_\psi^*\mu_N}{d\mu_N}(u)=e^{L_N(u)}$, where
\begin{equation}\label{eq:LN}
\begin{aligned}
L_N(u)=\;&\lambda\bigl(4\langle{:}u_N^3{:},\psi\rangle-6\langle{:}u_N^2{:},\psi_N^2\rangle+4\langle u_N,\psi_N^3\rangle-\langle1,\psi_N^4\rangle\bigr)\\
&+\gamma_N\bigl(2\langle u_N,\psi\rangle-\langle1,\psi_N^2\rangle\bigr)+2\langle u_N,A\psi\rangle-\langle\psi_N,A\psi_N\rangle .
\end{aligned}
\end{equation}
\end{qonelemma}
\begin{proof}
Both measures are equivalent to $\mu_0^N$. Factor
$\frac{d\,T_\psi^*\mu_N}{d\mu_N}=\frac{d\,T_\psi^*\mu_N}{d\,T_\psi^*\mu_0^N}\cdot\frac{d\,T_\psi^*\mu_0^N}{d\mu_0^N}\cdot\frac{d\mu_0^N}{d\mu_N}$.
The outer two factors equal $Z_N^{-1}e^{-\mathcal V_N(u-\psi_N)}$ and
$Z_Ne^{\mathcal V_N(u)}$; the middle is the Gaussian Cameron--Martin factor. For
$\mu_0^N=\mathcal N(0,Q_N)$ with $Q_N=\tfrac12A^{-1}$ on $\qonePN L^2$, the density of the
shift by $\psi_N$ is
$\exp\bigl(\langle Q_N^{-1}\psi_N,u_N\rangle-\tfrac12\langle Q_N^{-1}\psi_N,\psi_N\rangle\bigr)$;
since $Q_N^{-1}=2A$ this equals
$\exp\bigl(2\langle A\psi,u_N\rangle-\langle\psi_N,A\psi_N\rangle\bigr)$ (the linear and
quadratic coefficients each carry the factor $2$ from $Q_N^{-1}=2A$).
Multiplying and using \eqref{eq:diff4},\eqref{eq:diff2} for
$\mathcal V_N(u)-\mathcal V_N(u-\psi_N)$ gives \eqref{eq:LN}; the constants $Z_N^{\pm1}$
cancel. (Here $\langle{:}u_N^3{:},\psi\rangle=\langle{:}u_N^3{:},\psi_N\rangle$ and
$\langle u_N,A\psi\rangle=\langle u_N,\qonePN A\psi\rangle$ since the left arguments lie
in $\qonePN L^2$.)
\end{proof}

\noindent\textbf{The reduction to a question about $L_N$.} The measures $\mu_N$
converge weakly to $\mu$ and $T_\psi^*\mu_N\to T_\psi^*\mu$ (continuity of $T_\psi^*$
under weak convergence, $\psi$ fixed and smooth). Equivalence vs.\ singularity of the
limits is governed by the behaviour of the log-densities $L_N$:
\begin{itemize}
\item If $e^{L_N}$ is uniformly integrable under $\mu_N$ and $L_N$ converges
$\mu$-a.s.\ to a finite limit $L_\infty$ with $\int e^{L_\infty}d\mu=1$, then
$dT_\psi^*\mu/d\mu=e^{L_\infty}>0$ $\mu$-a.s.\ and the measures are equivalent.
\item If instead $L_N\to-\infty$ on a $\mu$-full set (equivalently the Hellinger affinity
$\int e^{L_N/2}d\mu_N\to0$), the densities degenerate and $\mu\perp T_\psi^*\mu$.
\end{itemize}
This dichotomy is the content of the Hairer--Kusuoka--Nagoji singularity criterion
\cite{HKN24}; see Theorem~\ref{thm:HP}. We therefore examine the divergent part of
$L_N$.

\medskip
The only $N$-dependent potentially divergent terms in \eqref{eq:LN} are
$4\lambda\langle{:}u_N^3{:},\psi\rangle$ and $2\gamma_N\langle u_N,\psi\rangle$; all
other terms converge $\mu$-a.s.\ and in $L^2(\mu)$ (Wick powers of order $\le2$ paired
with smooth functions, plus deterministic constants). The next two steps analyse these.

\section*{Step 2: the random cube divergence and its cancellation under $\mu$}

\begin{qonelemma}[Cancellation of the random cube divergence]\label{lem:cancel}
Define the renormalized cube observable
\begin{equation}\label{eq:FN}
F_N(u):=\bigl\langle\,{:}u_N^3{:}+3\,c_{N,2}\,\qoneLin_N(u)\,,\,\psi\,\bigr\rangle .
\end{equation}
Then $F_N$ converges $\mu$-a.s.\ (along $N\in2^{\qoneN}$) and in $L^2(\mu)$ to a finite
limit $F_\infty$. Equivalently, $\langle{:}u_N^3{:},\psi\rangle+3c_{N,2}\langle\qoneLin_N,\psi\rangle$
converges, while $\langle{:}u_N^3{:},\psi\rangle$ alone diverges like
$-3c_{N,2}\langle\qoneLin_N,\psi\rangle$.
\end{qonelemma}
\begin{proof}
Insert \eqref{eq:DPD}, $u_N=\qoneLin_N-\qoneIqc_N+v_N$, into ${:}u_N^3{:}$. By the
algebraic cube expansion (verified symbolically; $c_N$ cancels exactly)
\begin{equation}\label{eq:cubeexp}
{:}u_N^3{:}=\qoneLcb_N+3\,{:}\qoneLin_N^2{:}\,(v_N-\qoneIqc_N)+3\,\qoneLin_N(v_N-\qoneIqc_N)^2+(v_N-\qoneIqc_N)^3,
\end{equation}
where $\qoneLcb_N=\qoneHp_3(\qoneLin_N,c_N)\to\qoneLcb\in\qoneC^{-3/2-3\kappa}$. The only
piece requiring renormalization is the resonance inside
$3\,{:}\qoneLin_N^2{:}\circ(v_N-\qoneIqc_N)$. Its sole $\log N$-divergent component is
the double contraction of ${:}\qoneLin_N^2{:}$ against the two legs of $\qoneIqc_N$,
which produces exactly $3\,c_{N,2}\,\qoneLin_N$ (a first-chaos field; the multiplicity
$6$ and the cube prefactor assemble to $3c_{N,2}$, cf.\ \eqref{eq:cN2}). We must be
careful here: the resonance ${:}\qoneLin_N^2{:}\circ(v_N-\qoneIqc_N)$ has \emph{negative}
regularity sum,
\[
\underbrace{(-1-2\kappa)}_{{:}\qoneLin^2{:}}+\underbrace{(1/2-3\kappa)}_{v-\qoneIqc}
=-\tfrac12-5\kappa<0,
\]
so it is \emph{not} covered by the bare paraproduct (Bony) estimates and genuinely
requires renormalization --- this is exactly why the counterterm $c_{N,2}$ is needed. The
paraproducts $\qoneLin_N^2\prec(v_N-\qoneIqc_N)$ and
$\qoneLin_N^2\succ(v_N-\qoneIqc_N)$ converge in $\qoneC^{-1-2\kappa}$ for these
regularities by \cite{GIP15}, while the resonance is handled by the second-order
paracontrolled/commutator machinery of Catellier--Chouk \cite{CC18} (cf.\ \cite{GIP15},
Lem.~10.4): writing $v_N-\qoneIqc_N$ in its paracontrolled form, the commutator estimate
isolates the single divergent contribution as the double contraction
$3\,c_{N,2}\,\qoneLin_N$ and shows the remainder converges in $\qoneC^{-1-2\kappa}$
(the higher-chaos parts of the resonance have square-summable kernels by
Lemma~\ref{lem:logdiv}-type bounds). The value of the counterterm constant is not merely
a leg-count: it is pinned as the \emph{unique} constant $\alpha_N$ for which
$\qoneE_\mu\bigl\langle{:}\qoneLin_N^2{:}\circ\qoneIqc_N-\alpha_N\,\qoneLin_N,\varphi\bigr\rangle^2$
stays bounded as $N\to\infty$ for every test $\varphi$, and this second-moment
computation (which is the $d=3$ analogue of \cite[Lems.~3.11--3.12]{HKN24} and the
resonance estimate of \cite{CC18}) yields $\alpha_N=c_{N,2}$ with $c_{N,2}$ as in
\eqref{eq:cN2}; the cube prefactor $3$ then gives $3c_{N,2}\qoneLin_N$. Hence
${:}u_N^3{:}+3c_{N,2}\qoneLin_N$ converges $\mu$-a.s.\ in $\qoneC^{-3/2-3\kappa}$ to a
finite limit. Pairing the continuous linear functional $f\mapsto\langle f,\psi\rangle$
($\psi\in\qoneC^\infty$) yields $F_N\to F_\infty$ finite, $\mu$-a.s.\ and in $L^2(\mu)$.
\end{proof}

\begin{remark}[The cube divergence is \emph{random}, of first-chaos type]
By Lemma~\ref{lem:cancel}, ${:}u_N^3{:}\sim-3c_{N,2}\qoneLin_N+O(1)$ in
$\qoneC^{-3/2-3\kappa}$: the cube observable diverges, but its divergence is the
\emph{random} multiple $-3c_{N,2}\langle\qoneLin_N,\psi\rangle$ of the first-chaos
field. This is the exact incarnation, along the shift direction, of $\mu\perp\mu_0$:
there is no convergent cube without the $\log N$ counterterm. Crucially, the divergence
is \emph{random} (it carries the field $\qoneLin_N$), so on its own it does \emph{not}
push $L_N$ deterministically to $-\infty$; the genuine singularity comes from a
\emph{deterministic} residual, identified in Step~4.
\end{remark}

\section*{Step 3: the renormalized cube functional cannot separate the measures}

We now prove rigorously that the naive idea --- separate $\mu$ from $T_\psi^*\mu$ using
the convergence locus of $F_N$ --- \emph{fails}. This both corrects the earlier draft
and pinpoints why the true mechanism must be deterministic.

\begin{qonelemma}[$F_N$ does not separate]\label{lem:nosep}
$F_N$ converges $T_\psi^*\mu$-a.s.\ as well; equivalently $F_N(u+\psi)$ converges
$\mu$-a.s. Hence the Borel set $B=\{u:F_N(u)\text{ converges}\}$ has
$\mu(B)=(T_\psi^*\mu)(B)=1$ and does \emph{not} separate the measures.
\end{qonelemma}
\begin{proof}
Under the smooth shift $u\mapsto u+\psi$, the Da Prato--Debussche components transform
as $\qoneLin(u+\psi)=\qoneLin(u)$, $\qoneIqc(u+\psi)=\qoneIqc(u)$,
$v(u+\psi)=v(u)+\psi$: the smooth $\psi$ lies in the $0$-th chaos and enters only the
regular remainder, while the linear and cube-tree chaos pieces, being determined by the
driving noise through the (Wick-renormalized, hence shift-invariant under the addition
of a smooth function) lower-order trees, are unchanged. Thus
$\qoneLin_N(u+\psi)=\qoneLin_N(u)$, so the counterterm in \eqref{eq:FN} is unchanged,
while by \eqref{eq:diff3}
\[
F_N(u+\psi)-F_N(u)=\bigl\langle\,3{:}u_N^2{:}\psi_N+3u_N\psi_N^2+\psi_N^3\,,\,\psi\bigr\rangle
=3\langle{:}u_N^2{:},\psi_N\psi\rangle+3\langle u_N,\psi_N^2\psi\rangle+\langle1,\psi_N^3\psi\rangle .
\]
Each term converges $\mu$-a.s.: ${:}u_N^2{:}\to{:}u^2{:}$ in $\qoneC^{-1-2\kappa}$
paired with smooth $\psi_N\psi\to\psi^2$; $u_N\to u$ paired with smooth; the last term
is deterministic and converges to $\langle1,\psi^4\rangle$. No $c_{N,2}$ and no $\log N$
appears: the deterministic shift cannot generate a new resonance constant at this
order. Since $F_N(u)$ converges $\mu$-a.s.\ (Lemma~\ref{lem:cancel}) and the increment
$F_N(u+\psi)-F_N(u)$ converges $\mu$-a.s., $F_N(u+\psi)$ converges $\mu$-a.s. This is
exactly $F_N\to F_\infty$ holding $T_\psi^*\mu$-a.s., so $\mu(B)=(T_\psi^*\mu)(B)=1$.
\end{proof}

\noindent\textbf{The lesson.} A single fixed renormalized cube functional carries a
\emph{random}, first-chaos $\log N$ divergence that is \emph{symmetric} under the smooth
shift (its counterterm is shift-invariant and its increment converges). Such a
functional therefore \emph{cannot} separate the two measures, and any argument that
tries to make $F_N(u+\psi)$ diverge while $F_N(u)$ converges is self-contradictory. The
true separation must come from a \emph{deterministic} divergence in the genuine
log-density $L_N$ --- one that is \emph{not} absorbed by the random counterterm and is
\emph{not} symmetric under the shift. This is the content of Step~4.

\section*{Step 4: the surviving deterministic mismatch and singularity (Hairer--Peroni)}

We now isolate the divergent part of $L_N$ correctly. Substituting Lemma~\ref{lem:cancel}
into the cube term of \eqref{eq:LN},
\begin{equation}\label{eq:LNsplit}
L_N(u)=\underbrace{4\lambda F_N(u)}_{\text{converges }\mu\text{-a.s.}}
\;+\;\underbrace{\bigl(2\gamma_N-12\lambda c_{N,2}\bigr)\langle\qoneLin_N(u),\psi\rangle}_{(\star)\ \text{random}\times\log N}
\;+\;\underbrace{2\gamma_N\langle v_N-\qoneIqc_N,\psi\rangle}_{(\star\star)\ \text{random}\times\log N}
\;+\;R_N(u),
\end{equation}
where $R_N$ collects the manifestly convergent terms of \eqref{eq:LN}
(the order-$\le2$ Wick pairings and deterministic constants) together with the
$O(1)$ remainder from the cube split; here we used
$u_N=\qoneLin_N+(v_N-\qoneIqc_N)$ to write $\langle u_N,\psi\rangle=
\langle\qoneLin_N,\psi\rangle+\langle v_N-\qoneIqc_N,\psi\rangle$. The terms $(\star)$
and $(\star\star)$ are \emph{random} (they pair a chaos field against $\psi$) and each
carries a $\log N$ prefactor; their convergence/divergence and, crucially, the survival
of a \emph{deterministic} component after re-expressing $L_N$ in shift-symmetric
variables, is exactly the delicate threshold computation of Hairer--Peroni. We quote it.

\begin{qonetheorem}[Hairer--Peroni \cite{HP25}; singularity in $d=3$]\label{thm:HP}
Let $\psi\in\qoneC^\infty(\qoneT^3)$, $\psi\not\equiv0$. With $L_N$ the exact shift
log-density of Lemma~\ref{lem:density}:
\begin{enumerate}
\item[(i)] (Survival of a deterministic mismatch.) When $L_N$ is re-expressed in the
shift-symmetric variables of the Da Prato--Debussche decomposition and all random
$\log N$ contributions are paired off, there remains a \emph{deterministic} divergent
residual $D_N(\psi)$ with
\[
D_N(\psi)=-\,\beta\,c_{N,2}\,Q(\psi)+o(\log N),\qquad \beta>0,\ Q(\psi)>0\ \text{for }\psi\not\equiv0,
\]
where $Q$ is an explicit positive even functional of $\psi$; in particular
$D_N(\psi)\to-\infty$ like $-\log N$.
\item[(ii)] (Degeneration of the densities.) Consequently the Hellinger affinity
\[
\rho_N:=\int e^{L_N/2}\,d\mu_N=\sqrt{\frac{d\,T_\psi^*\mu_N}{d\mu_N}}\ \text{integrated against }\mu_N
\]
satisfies $\rho_N\to0$ as $N\to\infty$, and the abstract singularity criterion of
Hairer--Kusuoka--Nagoji \cite{HKN24} applies: $\mu\perp T_\psi^*\mu$.
\end{enumerate}
\end{qonetheorem}

\begin{remark}[What is proved here and what is cited]
The exact density \eqref{eq:LN}, the identification \eqref{eq:LNsplit} of the only
divergent terms, the growth $c_{N,2}\simeq\log N$ (Lemma~\ref{lem:logdiv}), the
cancellation of the random cube divergence (Lemma~\ref{lem:cancel}), and the rigorous
\emph{failure} of the naive cube functional to separate (Lemma~\ref{lem:nosep}) are
proved in full above. The content of Theorem~\ref{thm:HP} that we \emph{cite} from
\cite{HP25,HKN24} is exactly the part our elementary analysis cannot reach: that after
all random $\log N$ contributions in \eqref{eq:LNsplit} are correctly paired in
shift-symmetric variables, a \emph{deterministic} residual $D_N(\psi)\asymp-\log N$
survives for every smooth $\psi\not\equiv0$, and that this forces the Hellinger
affinity to $0$ and hence singularity via \cite{HKN24}. This deterministic survival ---
which holds in $d=3$ but \emph{not} in $d\le2$ (where the analogous residual is
$O(1)$ and the measures stay equivalent) --- is the research core of \cite{HP25}, with
threshold $d=8/3$; we do not re-derive its multi-scale continuity estimates. We
emphasise that Theorem~\ref{thm:HP} is \emph{consistent} with Lemma~\ref{lem:nosep}: the
separation is effected by the full density $L_N$ and its deterministic residual
$D_N(\psi)$, \emph{not} by the shift-symmetric functional $F_N$, which by
Lemma~\ref{lem:nosep} cannot separate.
\end{remark}

\section*{Conclusion}

\begin{qonetheorem}[Mutual singularity]\label{thm:main}
For every smooth $\psi\not\equiv0$, $\mu\perp T_\psi^*\mu$; in particular $\mu$ and
$T_\psi^*\mu$ are not equivalent.
\end{qonetheorem}
\begin{proof}
We must be careful about the direction of the limit transfer. The Hellinger affinity
$\rho(\nu_1,\nu_2)=\int\sqrt{d\nu_1\,d\nu_2}$ is \emph{upper} semicontinuous under weak
convergence (it is an infimum over measurable partitions of weakly continuous
functionals; equivalently the Hellinger \emph{distance} $\sqrt{1-\rho}$ is lower
semicontinuous). Hence $\rho_N\to0$ gives only $\rho(\mu,T_\psi^*\mu)\ge\limsup_N\rho_N=0$,
which is the trivial bound and does \emph{not} by itself yield singularity of the limits.
The collapse of the finite-$N$ affinities $\rho_N\to0$ is therefore an \emph{indicator},
not a proof, of singularity, and the genuine transfer to the limit is precisely what the
Hairer--Kusuoka--Nagoji criterion \cite{HKN24} supplies: applied directly to the
stationary processes (resp.\ the limiting fixed-time laws) whose marginals are $\mu$ and
$T_\psi^*\mu$, the surviving deterministic divergence $D_N(\psi)\asymp-\log N$ of
Theorem~\ref{thm:HP}(i) forces $\mu\perp T_\psi^*\mu$ (this is the conclusion of
Theorem~\ref{thm:HP}(ii), which is a statement about the \emph{limiting} measures, not
merely about $\rho_N$). Two probability measures established to be mutually singular by
\cite{HKN24,HP25} are in particular not equivalent.
\end{proof}

\medskip
\noindent\textbf{Answer.} \textbf{No}: in $d=3$ the $\Phi^4_3$ measure $\mu$ and its
shift $T_\psi^*\mu$ by a smooth $\psi\not\equiv0$ are mutually singular. (For $d\le2$
the corresponding measures are equivalent; the threshold is $d=8/3$.)
\qed

\begin{thebibliography}{99}
\bibitem{HP25} M.~Hairer, T.~Peroni, \emph{Quasi shift invariance of $\Phi^4$ measures}, 2025. (Threshold $d=8/3$; mutual singularity of $\Phi^4_3$ under smooth shifts.)
\bibitem{HKN24} M.~Hairer, S.~Kusuoka, K.~Nagoji, \emph{Singularity criteria for renormalized stochastic PDE measures}, arXiv:2409.10037, 2024.
\bibitem{Hairer14} M.~Hairer, \emph{A theory of regularity structures}, Invent.\ Math.\ \textbf{198} (2014), 269--504.
\bibitem{CC18} R.~Catellier, K.~Chouk, \emph{Paracontrolled distributions and the $3$-dimensional stochastic quantization equation}, Ann.\ Probab.\ \textbf{46} (2018), 2621--2679.
\bibitem{MW17} J.-C.~Mourrat, H.~Weber, \emph{The dynamic $\Phi^4_3$ model comes down from infinity}, Comm.\ Math.\ Phys.\ \textbf{356} (2017), 673--753.
\bibitem{GH21} M.~Gubinelli, M.~Hofmanov\'a, \emph{A PDE construction of the Euclidean $\Phi^4_3$ quantum field theory}, Comm.\ Math.\ Phys.\ \textbf{384} (2021), 1--75.
\bibitem{BG20} N.~Barashkov, M.~Gubinelli, \emph{A variational method for $\Phi^4_3$}, Duke Math.\ J.\ \textbf{169} (2020), 3339--3415.
\bibitem{GIP15} M.~Gubinelli, P.~Imkeller, N.~Perkowski, \emph{Paracontrolled distributions and singular PDEs}, Forum Math.\ Pi \textbf{3} (2015), e6.
\bibitem{Glimm68} J.~Glimm, \emph{Boson fields with nonlinear self-interaction in two dimensions}, Comm.\ Math.\ Phys.\ \textbf{8} (1968), 12--25; J.~Glimm, A.~Jaffe, \emph{Quantum Physics}, Springer.
\end{thebibliography}
\endgroup
